\documentclass[11pt,a4paper]{article}
\usepackage[margin=23mm]{geometry}
\usepackage{amsmath,amssymb,amsthm,array,booktabs,microtype}
\usepackage[hidelinks]{hyperref}
\setlength{\parindent}{0pt}
\setlength{\parskip}{6pt}
\newtheorem{theorem}{Theorem}
\newcommand{\C}{\mathbb C}
\newcommand{\V}{\mathcal V}
\newcommand{\Simp}{\operatorname{Simple}}
\newcommand{\Sub}{\operatorname{Sub}}
\begin{document}
\begin{center}
{\Large\bfseries A multifusion category with unit length three}\par
\smallskip
Problem 00000005688 \quad lizaixi01 \quad 24 October 2026
\end{center}

\textbf{Precise scope.}
The printed conjecture includes the universal clause ``the length of the unit is 2''.
We refute this clause by a genuine finite semisimple, complex-linear, rigid monoidal
category with non-simple unit. The category is decomposable. The source imposes no
indecomposability assumption; that property is separate from the definition of
multifusion in EGNO, Definition 4.1.1 [1, p.~65]. We make no claim about the other,
unspecified dimension-table, global-dimension or local-unit assertions.

\textbf{The category and its structure.}
Let $I=\{0,1,2\}$ and let $\V=\mathrm{Vec}_{\C}^{\mathrm{fd}}$. Set
\[
  \mathcal A=\operatorname{Fun}(\operatorname{Discrete}(I),\V).
\]
Thus an object is an arbitrary triple $(V_0,V_1,V_2)$ of finite-dimensional complex
vector spaces, and every triple of linear maps is a morphism. This is the full
category of such triples, rather than a list of selected objects. Tensor product,
unit and dual are componentwise:
\[
 (V\otimes W)_i=V_i\otimes_{\C}W_i,
 \qquad \mathbf1=(\C,\C,\C),\qquad (V^*)_i=V_i^*.
\]
The associators, unitors, evaluations and coevaluations are the actual vector-space
maps. Their coherence and duality identities are supplied by the native rigid
monoidal category interfaces. Kernels and cokernels are componentwise, so
$\mathcal A$ is abelian and $\C$-linear. Tensor is additive and $\C$-linear in both
morphism variables. Each Hom space is a finite product of finite-dimensional
linear-map spaces, hence finite-dimensional. The scalar field is the actual
algebraically closed field $\C$.

Write $S_i$ for the object with $\C$ in position $i$ and zero spaces elsewhere.
For every $V$, choose actual finite bases $b_{i,u}$ of $V_i$. The maps that insert
$b_{i,u}$ and extract its coordinate satisfy
\[
 \pi_{j,v}\,\iota_{i,u}=\delta_{(i,u),(j,v)}\operatorname{id}_{S_i},
 \qquad \sum_{i,u}\iota_{i,u}\,\pi_{i,u}=\operatorname{id}_V.
\]
They give an actual finite biproduct bicone, proved to be a bilimit:
\[
 V\ \cong\ \bigoplus_{i\in I}\ \bigoplus_{u<\dim_{\C}V_i} S_i.
\]
The classification of all subobjects on the next page proves that each $S_i$ is
simple. A nonzero basis insertion into any simple $V$ is an isomorphism by Schur's
lemma. Consequently these are all simple isomorphism classes. They are pairwise
nonisomorphic by evaluating at a coordinate. The genuine quotient of all simple
objects by existence of an isomorphism therefore has cardinality three.

\newpage
\textbf{All subobjects, finite length and the actual unit.}
For every object $V$ there is an order isomorphism
\begin{equation}\label{eq:subobjects}
 \Sub_{\mathcal A}(V)\ \cong\
 \prod_{i\in I}\Sub_{\C}(V_i).
\end{equation}
It sends a categorical monomorphism to the ranges of its three actual component
linear maps. The inverse takes a triple of subspaces to its inclusion into $V$.
The component map from the domain of a monomorphism to its range is a linear
isomorphism; these three isomorphisms form the natural isomorphism that verifies
the inverse law. Inclusion of ranges is equivalent to categorical factorization.
Thus (\ref{eq:subobjects}) describes every subobject, not only coordinate examples.

For $S_i$, the only choices are $0$ or $\C$ at coordinate $i$, while the other
coordinates have only their zero subspace. Its subobject order is therefore
simple, giving native categorical simplicity. For arbitrary $V$, the integer
\[
 r(N)=\sum_{i\in I}\dim_{\C}N_i
\]
strictly increases on every strict inclusion of subobjects and is at most
$\sum_i\dim_{\C}V_i$. Every strict chain has bounded length; the native
\texttt{FiniteDimensionalOrder} predicate for the actual subobject order follows.
Together with the finite simple biproducts above, this proves finite semisimplicity
on all objects. Epimorphisms are componentwise surjections. Linear lifting through
these surjections proves that every object is projective, and hence the category
has enough projectives.

The actual bases of the three unit components each have one element. Reindexing
their insertion and projection maps gives a bilimit with precisely the family
$(S_0,S_1,S_2)$, and therefore
\[
                 \mathbf1\ \cong\ S_0\oplus S_1\oplus S_2.
\]
The unit is non-simple: every simple object has total vector-space dimension one,
whereas the actual unit has total dimension three, an isomorphism invariant.

\textbf{Composition length, rather than only vector-space dimension.}
At the unit, (\ref{eq:subobjects}) gives an order isomorphism with $\{0,1\}^3$,
because each actual subspace of $\C$ is zero or the whole line. The chain
\[
 (0,0,0)\ \lessdot\ (1,0,0)\ \lessdot\ (1,1,0)\ \lessdot\ (1,1,1)
\]
consists of covers: exactly one coordinate changes, and there is no intermediate
subspace. Transport through the order isomorphism constructs a native
\texttt{CompositionSeries} of $\Sub_{\mathcal A}(\mathbf1)$ with bottom and top
endpoints and length three. The subobject lattice is modular by
(\ref{eq:subobjects}) and the modularity of subspace lattices. The native
Jordan--H\"older theorem proves that \emph{every} complete composition series of
this actual unit has length three. Endpoints are imposed explicitly, since the
native series type also allows intervals. In particular no complete composition
series of length two exists. This is the usual composition length of an object
in a finite-length abelian category, and refutes the printed universal clause.

\newpage
\textbf{Formal implementation and reproducibility.}
The accompanying Lean project uses the genuine native
\texttt{FGModuleCat}\,$\C$ at universe zero, and its functor category from the
discrete category on \texttt{Fin 3}. No category properties, dimensions or length
values are assigned as hypotheses. The proof constructs the maps and proves the
following bridges.

\begin{center}
\small
\begin{tabular}{@{}p{0.44\textwidth}p{0.49\textwidth}@{}}
\toprule
Lean declaration & Mathematical role\\
\midrule
\texttt{basisBilimit} & All-object finite basis biproducts.\\
\texttt{subobjectSpaces} & Complete subobject--subspace order isomorphism.\\
\texttt{coordSimple} & Native categorical simplicity of each coordinate.\\
\texttt{all\_simples\_are\_coordinates} & Classification on all simple objects.\\
\texttt{three\_simple\_classes} & Cardinality of the actual isomorphism quotient.\\
\texttt{all\_objects\_finite\_length} & Bounded strict chains in actual subobjects.\\
\texttt{everyObjectProjective} & Actual lifting through all epimorphisms.\\
\texttt{unitBiproductIso} & Isomorphism from the actual tensor unit.\\
\texttt{unit\_composition\_length} & All complete native unit series have length three.\\
\texttt{no\_length\_two\_unit} & Exclusion of every complete length-two series.\\
\bottomrule
\end{tabular}
\end{center}

The final theorem \texttt{printed\_unit\_length\_two\_false} collects the actual
linear, abelian, rigid, bilinear, finite-Hom, finite-semisimple and unit-length
properties. The separate native projectivity instances provide projective
presentations. The complete original bilingual problem is retained in
\texttt{SOURCE.md}; the report proves only its explicit unit-length conjunct.

The project pins Lean 4.33.0 and the exact Mathlib revision recorded in its
manifest. In the \texttt{lean} directory, run the four independent commands:
\begin{verbatim}
lean --version
lake build +Main
lake env lean -DwarningAsError=true Check.lean
lake env leanchecker --fresh Main
\end{verbatim}
Alternatively run \texttt{lean --version} followed by \texttt{bash verify.sh};
the script runs the remaining three commands. The audit enumerates every
declaration owned by \texttt{TLMC5688}, including generated auxiliaries, and calls
ordinary \texttt{\#print axioms} for each. The only permitted foundational axioms
are \texttt{propext}, \texttt{Classical.choice} and \texttt{Quot.sound}.
No additional axiom, unfinished proof or native decision shortcut is used.
Execution times in the evidence are production diagnostics, not controlled
research-efficiency measurements.

\textbf{Definition boundary and reference.}
The category is the componentwise product of three vector-space categories. It
is not the matrix category $\mathrm{Mat}_3(\mathrm{Vec})$ of EGNO Example 4.1.3.
Reference [1] is used for the standard terminology and the separate status of
indecomposability; it supplies no unproved identification or numerical result to
Lean. All category, subobject, simplicity and length arguments used here are
proved in the accompanying project.

\begin{thebibliography}{1}
\bibitem{EGNO} P.~Etingof, S.~Gelaki, D.~Nikshych and V.~Ostrik,
\emph{Tensor Categories}, Mathematical Surveys and Monographs 205,
American Mathematical Society, 2015. Author final version,
Definition 4.1.1 and Examples 4.1.2--4.1.3, printed p.~65.
\url{https://math.mit.edu/~etingof/egnobookfinal.pdf}
\end{thebibliography}
\end{document}
